% =============================================================================
%  Part IX -- Discussion and conclusion
% =============================================================================

\section{What has been shown}\label{sec:shown}

Starting from events and their order, Causal Physics reaches the following, each step tested or derived:
\begin{enumerate}
\item The substrate is a Lorentz-invariant causal set, and its geometry can be read blindly from order and
  number.
\item Matter is a flow that carries its four-velocity; localisation requires this memory, and the memory is
  inertia. The flows have a conserved energy--momentum tensor.
\item The order responds to matter, and the way it responds fixes $\gamma$: a fixed order gives
  $\gamma=-1$, narrowing cones give $\gamma=1/d_s$, conservation of throughput gives $\gamma=2/(d_s-1)$,
  which is $1$ in three spatial dimensions.
\item In three spatial dimensions the throughput principle is the coupling of linearised general relativity
  to $T_{ab}$, and a substrate rule realising it makes pressure and motion gravitate and drags the cones.
\item That rule is not an independent postulate: with the proper time of a flow counted by the order, it is
  the linear response of an order whose action has the Einstein--Hilbert stiffness.
\item The discrete Benincasa--Dowker action has that stiffness: $6.0003$ against $6$ in the static
  conformal sector, $7.983$ against $8$ in the mixed sector and $2.002$ against $2$ in the optical
  sector, computed with nothing from general relativity put in; $\gamma=1$ follows from it. The discreteness
  adds a non-analytic correction and non-local terms that grow with the age of the field, and beyond the
  leading order the sectors differ.
\end{enumerate}

\section{Limitations}\label{sec:limitations}

\begin{itemize}
\item \textbf{Expected values, not dynamics.} The action of the order is evaluated as an expectation over
  sprinklings of a given geometry. That a dynamics of the order selects such causal sets, the manifold-like
  phase, is assumed, as in causal set theory generally.
\item \textbf{Finite ages.} The coefficients are computed for ages up to a few hundred $\ell_k$; the age term of
  the optical sector grows faster than linearly there, and its asymptotic form is not known. Beyond the
  leading order the sectors are not in covariant proportion; whether this comes from the age cut or from the
  action is open.
\item \textbf{The substrate rule as an input in the tests.} The tests of the coupling (F7.1c, F7.4a) build
  the order according to the rule. They show what the rule implies, not that nature uses it; the
  derivation from the action supplies the reason.
\item \textbf{Matter as a self-bound object.} Flows are localised by their memory of direction, but a
  particle that holds itself together without an external source has not been obtained.
\item \textbf{The quantum.} The theory is classical. Superposition and measurement have an
  interpretation, not a dynamics.
\item \textbf{The age term.} Its coefficient is established; its physical reading, what $T$ is, is not, and
  it may depend on how the past is cut. The reading decides whether short-range gravity tests probe $\ell_k$ now
  (Sec.~\ref{sec:window}).
\item \textbf{The scale $\ell_k$.} Existing bounds on non-local fields give $\ell_k\lesssim10^{-19}$\,m. Whether
  they apply to the gravitational sector depends on the matter sector, which is open.
\item \textbf{Finite numerics.} Laboratory results come from finite regions and finite numbers of
  sprinklings, with the limits stated in each test.
\end{itemize}

\section{Outlook}\label{sec:outlook}

The road from here to predictions is short and concrete: the asymptotic age term of the optical sector,
then the static field of a point mass from the three coefficients, compared with short-range gravity data,
and the physical meaning of the age $T$. In parallel, the thermodynamic route is to turn the completion by general relativity from
an adopted principle into a derived one, and the matter sector is to find the self-bound flow.

\section{Conclusion}\label{sec:conclusion}

Gravity in Causal Physics is the response of a causal order to the matter that flows through it. The
order slows its clocks, narrows its cones and adds events so that its throughput is kept, and in three
spatial dimensions this is exactly the coupling of general relativity. The reason is the action of the
order itself: counting the intervals of a causal set gives the Einstein--Hilbert stiffness, and the
counting also says where it stops being exact. That is where the predictions of the theory are.
